% !TEX root = ProjectDescription.tex
\section{Overview and Objectives}

Embedded systems are ubiquitous computing devices deployed across critical infrastructure, consumer electronics, medical devices, industrial control systems, and IoT applications.
These systems range from microcontroller units (MCUs) that power sensors and edge devices to more capable embedded processors in automotive systems, smart home devices, and industrial controllers.
They are essential to everyday life, with billions of embedded devices deployed globally and projected to grow to trillions by 2035~\cite{Philip2021route}.
While the benefits of these systems are unparalleled, unfortunately, they are susceptible to cyberattacks that are occurring at unprecedented levels, and which often have severe consequences ranging from loss of life~\cite{FDA2019} to homeland security breaches~\cite{stellios2018survey}.
To ensure our embedded and IoT infrastructure is built on a trustworthy and secure foundation, the ``IoT Cybersecurity Improvement Act of 2019" (S.734 in the Senate)~\cite{s734} highlights the need for urgent actions on embedded systems security.

In response to these persistent software vulnerabilities, the cybersecurity community and government agencies have pursued multiple complementary strategies to mitigate memory corruption attacks on embedded systems.
One major effort pursues defensive techniques that continue to use C and C++ while deploying compiler-based and runtime mitigations—including stack canaries, address space layout randomization (ASLR), data execution prevention (DEP), control-flow integrity (CFI), and software fault isolation (SFI)—to reduce the exploitability of memory corruption bugs without requiring language rewrites.
A parallel and increasingly strategic effort involves transitioning to memory-safe programming languages: the White House and federal agencies have invested significant effort in promoting memory-safe languages as a defense against memory corruption vulnerabilities~\cite{whitehouse2024memory}, with major initiatives including DARPA's memory-safety focused research programs~\cite{darpa2024ssith} underscoring the importance of moving embedded system development away from unsafe languages like C and toward memory-safe alternatives such as Rust.
Both approaches have made measurable progress: as we showed in our CCS 2025 paper analyzing MITRE's Embedded Capture-The-Flag (eCTF) competitions~\cite{ma2025ectf}, both the deployment of defensive techniques in C/C++ and the adoption of Rust are increasing, with memory-corruption-based vulnerabilities declining significantly across competition rounds.

Yet both memory-safe languages and hardened C/C++ share a fundamental limitation: they only protect the software layer.
Embedded systems operate under constraints that expose them to a different class of attacks altogether: they need hardware access, run on tight resource budgets, and execute on processors with minimal built-in security.
When software defenses mature, attackers simply shift their approach—they move to hardware-level attacks that exist below the software layer entirely.
In recent eCTF competitions (as documented in our CCS 2025 analysis~\cite{ma2025ectf}), we're seeing exactly this play out: multiple teams are successfully using fault injection to compromise embedded systems, whether those systems are running hardened C/C++ code or Rust.
This means neither approach solves the problem—fault injection is orthogonal to memory safety, and it's a threat that remains largely unaddressed in embedded systems security research.

\subsection{Proposed Research}

\begin{figure}[t]
\centering
\resizebox{\textwidth}{!}{%
\begin{tikzpicture}[
  font=\footnotesize,
  >={Stealth[length=2.2mm]},
  layer/.style={draw=blue!55, fill=blue!8, rounded corners, align=center, minimum width=3cm, minimum height=1.1cm},
  glitchbox/.style={draw=orange!80!black, fill=orange!18, rounded corners, align=center, font=\scriptsize},
  thrust/.style={draw=black!55, fill=black!4, rounded corners, align=center, minimum height=1.6cm},
  pipe/.style={->, very thick, blue!50!black},
  ann/.style={->, dashed, black!55},
]
% --- implementation layers (pipeline), spread across the full column ---
\node[layer] (src)  at (2,0)    {Source\\code};
\node[layer] (comp) at (7.5,0)  {\textbf{Compiler}\\[1pt]{\scriptsize optimization, instruction}\\{\scriptsize selection, reg.\ allocation}};
\node[layer] (bin)  at (13,0)   {Binary on\\embedded device};
\draw[pipe] (src)  -- (comp);
\draw[pipe] (comp) -- (bin);
% --- to-be-addressed issue: fault injection ---
\node[glitchbox] (glitch) at (13,2.1) {Fault injection\\(voltage / clock / EM / laser)};
\draw[->, very thick, orange!80!black] (glitch) -- (bin);
% --- thrusts in one row ---
\node[thrust, text width=4.0cm] (t3) at (2,-2.7)   {\textbf{Thrust 3} (\mbox{\emph{Defend}})\\[2pt]{\scriptsize Cost-aware hardening + fault-aware code generation; provable resistance under a \emph{bounded} fault model}};
\node[thrust, text width=4.0cm] (t1) at (7.5,-2.7) {\textbf{Thrust 1} (\mbox{\emph{Understand}})\\[2pt]{\scriptsize Characterize \& quantify the compiler-enlarged glitching attack surface}};
\node[thrust, text width=4.0cm] (t2) at (13,-2.7)  {\textbf{Thrust 2} (\mbox{\emph{Detect}})\\[2pt]{\scriptsize Automatically find exploitable faults under a fault model}};
\draw[ann] (t1.north) -- (comp.south);
\draw[ann] (t2.north) -- (bin.south);
% --- Thrust 3 feedback loop into the compiler ---
\draw[->, thick, blue!60!black, rounded corners=6pt]
  (t3.west) -- (-0.7,-2.7) -- (-0.7,1.7) -- (7.5,1.7) -- (comp.north);
\node[blue!60!black, font=\scriptsize] at (3.0,1.95) {harden the compiler};
\end{tikzpicture}%
}
\caption{Overview of \sysname. The three-thrust agenda follows an \emph{understand $\rightarrow$ detect $\rightarrow$ defend} arc unified by the compiler. To-be-addressed issues are shown in {\color{orange!85!black}orange} and implementation layers in {\color{blue!60!black}blue}: the compiler enlarges the glitching attack surface (Thrust~1), which we discover automatically (Thrust~2) and then reduce by feeding cost-aware, verifiable hardening back into the compiler (Thrust~3).}
\label{fig:overview}
\end{figure}

As shown in Figure~\ref{fig:overview}, this proposal pursues a three-thrust agenda, \sysname, that follows an \emph{understand $\rightarrow$ detect $\rightarrow$ defend} arc unified by the compiler.
% TODO(glitching): state the overarching scientific goal in 1-2 bold sentences (cf. previous proposal), then the per-thrust paragraphs below.

\noindent\textbf{Thrust 1: \thrustone.}
% Compilers (optimization passes, instruction selection, register allocation) can introduce
% single points of failure and enlarge the fault-injection surface. Characterize and attribute this.

\noindent\textbf{Thrust 2: \thrusttwo.}
% A fully automated system to discover glitching vulnerabilities under a fault model.
% Framed as defensive bug-finding.

\noindent\textbf{Thrust 3: \thrustthree.}
% Close the loop: re-engineer the compiler to reduce the surface selectively (minimal overhead)
% and prove residual resistance under a bounded fault model. (We reduce, not eliminate -- a
% determined attacker with enough attempts can always inject some fault.)

\noindent\textbf{Novelty and relation to prior work.}
\textbf{Our key observation is that the compiler shapes a system's exposure to \emph{physical} fault injection in the same way it is known to shape its exposure to \emph{software} security bugs, yet only the software side has been studied.}
A closely related line of prior research shows that compilers can introduce software security weaknesses: optimizations delete memory-scrubbing operations~\cite{yang2017deadstore}, exploit undefined behavior to remove safety and security checks~\cite{wang2013undefined}, and more broadly open a ``correctness-security gap'' in which a transformation that preserves functional correctness nonetheless violates a security property~\cite{dsilva2015correctness}; a recent systematic study finds these compiler-introduced security bugs widespread across modern toolchains and optimization passes~\cite{xu2023silent}.
Our own work in the MITRE embedded CTF competition makes the pattern concrete: compiler optimizations silently removed \texttt{memset} calls intended to scrub AES keys from the stack, leaving sensitive data recoverable after it was meant to be erased~\cite{ma2025ectf}.
\sysname carries the same observation to physical fault injection: to our knowledge, no prior work characterizes, detects, or reduces the glitching attack surface that the compiler introduces, and \sysname is the first to treat the compiler as both the cause and the remedy of this physical-layer exposure.

\subsection{Relevance to SaTC}

This proposal directly serves the mission of the NSF SaTC~2.0 program (NSF~25-515), which ``aims to build trust in global cyber ecosystems'': ecosystems that, as the solicitation observes, have grown to involve ``hardware, software, networks, data, people, organizations, countries, and the physical world.''
\sysname addresses the seam between software and the physical world: fault-injection (glitching) attacks that violate the assumptions software, and the compilers that produce it, make about correct execution.
The effort falls squarely within the program's \emph{Computing and Communication Systems} scope, spanning hardware security, embedded systems, and trusted execution, and it also contributes to \emph{Foundations} through the formal, bounded-fault-model resistance guarantees developed in Thrust~3.
By characterizing (Thrust~1), automatically discovering (Thrust~2), and reducing (Thrust~3) compiler-introduced glitching vulnerabilities, \sysname strengthens the trustworthiness of the embedded and IoT devices that underpin transportation, healthcare, energy, and other critical infrastructure, advancing SaTC's goal of a resilient cyber ecosystem.
% TODO(glitching): if pursuing the optional TTE (Transition to Education) designation, add 1-2 sentences here on the integrated education plan.

\subsection{PI Qualifications}

The two PIs bring complementary expertise across the compiler and the hardware/software boundary that \sysname targets, on top of a broad record in systems and software security~\cite{zhang2019icore, zhao2011automatic, zhao2013using, tang2017nivanalyzer, jing2014morpheus, jing2014riskmon, jing2015towards, li2022understanding}.

\noindent\textbf{PI Zhao (Lead PI, Northeastern University)} is an Associate Professor at the Khoury College of Computer Sciences who has focused extensively on embedded system security and on the microarchitectural hardware/software boundary, the same layer at which glitching operates.
He designed CacheLight~\cite{gutierrez2018cachelight} and SmokeBomb~\cite{cho2020smokebomb} to mitigate cache side-channel attacks on the ARM architecture, and uncovered the Prime+Count cross-world covert channel on ARM TrustZone~\cite{cho2018prime}.
His group has crawled thousands of real-world embedded firmware samples to identify and analyze their vulnerabilities~\cite{nino24firware}, built trusted execution environments for embedded and IoT systems~\cite{zhao24tee, byotee24}, designed low-overhead control-flow attestation for microcontroller-based systems~\cite{armanuzzaman2025enola}, and works directly with electromagnetic emanations from embedded devices~\cite{Moradihaghighi2026fuzzemup}, the same physical channel exploited by EM fault injection.
Like fault injection, these attacks exploit physical and microarchitectural behavior rather than software bugs, giving PI Zhao direct experience with the threat model, measurement, and defenses this proposal targets.
His research has garnered several awards, including the \textbf{SACMAT Test-of-Time Award 2024} and \textbf{distinguished paper awards} at \textbf{USENIX Security 2019}, \textbf{AsiaCCS 2022}, and \textbf{CODASPY 2014}.
% TODO(glitching): verify award attribution now that the team changed -- the prior proposal listed these as team awards shared with Co-PI Kirda; NDSS 2011 (Kirda's EXPOSURE) has been removed.
He co-founded the USENIX Symposium on Vehicle Security and Privacy (VehicleSec) as modern vehicles are full of embedded systems, and has a track record of open-sourcing prototypes~\cite{ma2023dac, zhang2018secore, cho2020smokebomb, fengformal2021, lamp2019exsol} and transitioning them to patents and products~\cite{cho2020smokebomb, ahn2015methods, tu2017systems, isign}, e.g., in Samsung Tizen OS~\cite{Tizen}.

\noindent\textbf{PI Tan (University of Colorado Colorado Springs)} is an Assistant Professor who specializes in the security of embedded and ARM Cortex-M systems, with a focus on control-flow security and the architectural limits of deploying modern defenses on microcontroller-based systems.
Her work in this area includes a comprehensive bottom-up study of Arm Cortex-M security~\cite{tan24sok}, an analysis of the effectiveness of stack canaries on microcontroller systems~\cite{tan24canary}, and SHERLOC, a secure and holistic control-flow violation detector for embedded systems~\cite{tan2023sherloc}.
PI Tan will lead the hardware glitching testbed and on-device validation effort (Thrust~2). % TODO(glitching): confirm title/dept; add fault-injection-specific works if available.

\noindent\textbf{Joint efforts.}
The PIs have a sustained record of working together on embedded control-flow security, including attack-and-defense studies of return-to-non-secure and cross-state control-flow hijacking on Cortex-M TrustZone~\cite{ma2023dac, ma2025microft} and debugger-based asynchronous control-flow integrity for embedded systems~\cite{wang24insect}, giving them shared familiarity with the compiler and hardware/software boundary that glitching subverts.
Most relevant to \sysname, the PIs have jointly led teams in the MITRE Embedded Capture-the-Flag (eCTF) competition for six consecutive years and have published studies distilling its security lessons~\cite{ma2025ectf, ma2026wejust}, placing among the top teams \TODO{confirm years/placements}.
In these competitions, the team used power/voltage glitching and related fault-injection techniques to bypass secure boot, defeat firmware protections, and extract secrets from the embedded targets \TODO{add concrete examples}, giving the PIs an attacker's-eye view of the threats \sysname sets out to characterize, find, and reduce.
